\section{Diagonal Lerch velocities and their classical inverse series}
\label{sec:diagonal-inverse}

The harmonic-order continuation report \cite{HarmonicBaseline} defines,
for integers $n\geq 0$, $k\geq 1$, lattice spacing $h>0$, and
$0\leq\rho\leq1$,
\[
 f_{n,k}(a)=(-1)^{n+k}\frac{d^k}{da^k}
                  \left(\frac{\log^n a}{a}\right),\qquad
 \mathscr F_{n,k}^{\rho,h}(a)
       =\sum_{m\geq0}\rho^m f_{n,k}(a+mh).
\]
For fixed $n,k$, the derivative series converges locally uniformly on
$a>0$, including at $\rho=1$.  When $k=n-1$ and $n\geq2$, the smallest
positive elementary zero is $a=1$, it is simple, and
$f_{n,n-1}'(1)=-n!$.  Consequently it has an analytic local continuation
$a_n^{(h)}(\rho)$ near $\rho=0$.  Put $A=1+h>1$ and
\begin{equation}\label{eq:velocity-normalization}
 v_n(A)=-\frac1{n!}
   \left.\frac{d^{n-1}}{da^{n-1}}
                  \left(\frac{\log^n a}{a}\right)\right|_{a=A}.
\end{equation}
Implicit differentiation gives
$(a_n^{(h)})'(0)=v_n(A)$ for $n\geq2$.  The auxiliary coefficient
$v_1(A)=-(\log A)/A$ completes the generating series; no zero-branch
interpretation is asserted for this auxiliary index.

Let $\stirling{n}{j}$ denote the unsigned Stirling numbers of
the first kind and write $L=\log A$.  The baseline proves the exact
identities
\begin{align}
 v_n(A)&=A^{-n}\sum_{j=1}^n
           \stirling{n}{n-j+1}\frac{(-L)^j}{j!},
                  \label{eq:velocity-stirling}\\
 V_A(t)&=\sum_{n\geq1}v_n(A)t^n=L-U_A(t),\qquad
 e^{U_A(t)}=A+tU_A(t),\qquad U_A(0)=L.
                  \label{eq:velocity-inverse}
\end{align}
For clarity, the second formula is initially a statement about analytic
germs at zero.  It follows from Lagrange inversion \cite{Gessel2016} applied to
$X=A+t\log X$, with $U_A=\log X$; the analytic implicit function theorem
applies since the derivative with respect to $U$ at $(0,L)$ is $A\neq0$.

The polynomial $p_n(L)=A^n v_n(A)$ is also, up to the explicit sign
$L_n(L)=(-1)^n p_n(L)$, Cohen's Lambert polynomial
\cite[Proposition~4.8]{Cohen2021}.  This identification connects the
velocity problem to the real-root and extreme-root questions treated
in Sections~\ref{sec:velocity-zero-law} and~\ref{sec:cohen-edges}.

These coefficient polynomials have a classical antecedent that should be
made explicit.  Kalugin and Jeffrey \cite[\S2.1, Theorem~1]{KaluginJeffrey2012}
consider
\[
 1-e^{-v}+\sigma v-\sigma\lambda=0,\qquad v(0)=0.
\]
The exact change of variables
\begin{equation}\label{eq:classical-variable-map}
 \lambda=L,\qquad \sigma=-t/A,\qquad v=V_A(t)
\end{equation}
identifies their fixed-$\lambda$ series with \eqref{eq:velocity-inverse}.
Their double series, specialized by $\tau=\lambda\sigma$, gives exactly
\eqref{eq:velocity-stirling}.  Their convergence criterion also gives the
radius obtained below.  Thus neither the coefficient array nor the
radius formula is presented as a new discovery in the broader literature.
The purpose here is to give a direct barrier proof with full sheet
selection and a radial-slit continuation, thereby settle the precisely
stated accessibility conjecture in \cite{HarmonicBaseline}, and derive
its quantitative consequences for the Lerch zero velocities.

\section{An exact barrier selects the dominant conjugate pair}
\label{sec:dominant-pair}

For $A>1$ define $\theta\in(0,\pi)$ by
\begin{equation}\label{eq:theta-definition}
 1-\theta\cot\theta+
       \log\left(\frac{\theta}{\sin\theta}\right)=\log A,
\end{equation}
and set
\begin{equation}\label{eq:critical-parameters}
 a=1-\theta\cot\theta,\qquad u_*=a+i\theta,\qquad
 \tau=e^{u_*},\qquad R=e^a.
\end{equation}
The left side of \eqref{eq:theta-definition} increases from zero to
infinity, since its derivative is
\[
 \theta\csc^2\theta+\frac1\theta-2\cot\theta
 =\frac{(\theta-\sin\theta\cos\theta)^2+\sin^4\theta}
        {\theta\sin^2\theta}>0.
\]
Hence $\theta$ exists and is unique.  Moreover $a>0$, since
$\sin\theta-\theta\cos\theta>0$ on $(0,\pi)$, and
\begin{equation}\label{eq:critical-relations}
 \frac AR=\frac{\theta}{\sin\theta},\qquad
 (1-u_*)\tau=A.
\end{equation}
Equivalently, $u_*=1+W_0(-A/e+i0)$, using the upper boundary value of
the principal Lambert branch \cite{DLMFLambert}.  The angle definition
avoids a dependence on boundary-value conventions.

\begin{theorem}[Selected inverse sheet and its exact singularities]
\label{thm:selected-sheet}
For every $A>1$, the analytic germ $U_A$ in
\eqref{eq:velocity-inverse} has radius of convergence exactly $R$.
Its only singularities on $|t|=R$ are $\tau$ and $\overline\tau$.
There exists $\epsilon>0$ such that it has a single-valued holomorphic
continuation to
\begin{equation}\label{eq:slit-domain}
 \{t:|t|<R+\epsilon\}\setminus
 \left(\{s\tau:1\leq s<1+\epsilon/R\}
       \cup\{s\overline\tau:1\leq s<1+\epsilon/R\}\right).
\end{equation}
At the upper endpoint, writing $z=1-t/\tau$, the continued branch has a
convergent expansion
\begin{equation}\label{eq:selected-puiseux}
 U_A(t)=u_*+\chi z^{1/2}+b z+c_3z^{3/2}+b_2z^2+O(z^{5/2}),
\end{equation}
where
\begin{equation}\label{eq:puiseux-first-coefficients}
 \chi=\sqrt{-2u_*},\quad\Re\chi>0,\quad\Im\chi<0,
 \qquad b=\frac{u_*}{3}-1,
 \qquad c_3=\chi\left(\frac12-\frac{u_*}{18}-\frac1{4u_*}\right).
\end{equation}
Here $z^{1/2}$ is positive on the inward radial segment $t=s\tau$,
$s<1$.  The lower-endpoint expansion is the complex conjugate.
\end{theorem}

\begin{proof}
Write $F_t(u)=e^u-A-tu$.  We first construct a domain in the $u$-plane
that contains exactly one root for every $|t|<R$.  Put $d=1-a$ and
consider the rectangle
\[
 D_M=\{u=x+iy:0<x<M,\ |y|<\theta\},
\]
where $M>L$ will be chosen large.  On its upper horizontal side, set
$h=x-a$.  Equations \eqref{eq:critical-relations} give the exact
factorization
\begin{align}
 &|e^{x+i\theta}-A|^2-R^2(x^2+\theta^2)\notag\\
 &\quad=R^2\bigl((e^h-d)^2-(h+1-d)^2\bigr)\notag\\
 &\quad=R^2(e^h-1-h)(e^h+1+h-2d).
                         \label{eq:barrier-factorization}
\end{align}
The same identity holds on the lower side.  For $x\geq0$ we have
$h\geq-a$, and the second factor is increasing in $h$; its value at
$h=-a$ is $e^{-a}-1+a>0$.  The first factor is nonnegative and vanishes
only at $h=0$.  Therefore
\begin{equation}\label{eq:horizontal-bound}
 |e^u-A|\geq R|u|\quad(\Im u=\pm\theta,\ \Re u\geq0),
\end{equation}
with equality only at $u_*$ on the upper side and $\overline{u_*}$ on
the lower side.

On the left side, for $0<|y|\leq\theta$,
\begin{equation}\label{eq:vertical-bound-formula}
 \frac{|e^{iy}-A|^2}{y^2}
   =\frac{(A-1)^2}{y^2}
      +A\left(\frac{\sin(y/2)}{y/2}\right)^2.
\end{equation}
As a function of positive $y\leq\theta<\pi$, both terms decrease and
the first decreases strictly.  Its minimum occurs at $y=\theta$,
where \eqref{eq:barrier-factorization} is strictly positive because
$x=0\neq a$.  Thus $|e^{iy}-A|>R|y|$ on the entire left side; at
$y=0$ this is simply $A-1>0$.  Choose $M$ so large that
\[
 e^M-A>R\sqrt{M^2+\theta^2}.
\]
The right side then satisfies the same strict bound.

For $|t|<R$ the full boundary of $D_M$ consequently satisfies
$|tu|<|e^u-A|$.  Rouch\'e's theorem says that $F_t$ and $e^u-A$ have
the same number of zeros in $D_M$, counted with multiplicity.  The
latter has exactly the one simple zero $L$: the other zeros
$L+2\pi ik$ are excluded by $\theta<\pi$.  Hence $F_t$ also has exactly
one zero, and it is simple.  The implicit-function branches glue by
uniqueness, so $U_A$ is holomorphic for $|t|<R$, with
\begin{equation}\label{eq:root-strip}
 0<\Re U_A(t)<M,\qquad |\Im U_A(t)|<\theta.
\end{equation}

If $|t_0|=R$ and $t_0\notin\{\tau,\overline\tau\}$, then
$F_{s t_0}$ has no boundary zero for $0\leq s\leq1$.  For $s<1$ this
follows from the strict bounds.  At $s=1$ a boundary zero could occur
only at one of the two equality points in \eqref{eq:horizontal-bound},
and then \eqref{eq:critical-relations} would force $t_0$ to be the
corresponding excluded parameter.  The argument-principle count is
therefore one throughout the homotopy.  At $t_0$ there is precisely one
simple root inside $D_M$, and the implicit function theorem continues
$U_A$ through $t_0$.

It remains to identify the branch at a critical endpoint.  Set
$t=\tau(1-z)$ and $u=u_*+w$.  The exact local equation is
\begin{equation}\label{eq:local-implicit-equation}
 e^w-1-w+z(u_*+w)=0.
\end{equation}
Since $u_*\neq0$, its two local inverses are convergent power series
in $z^{1/2}$, beginning $w=\pm\chi z^{1/2}+O(z)$ with the choice
of $\chi$ in \eqref{eq:puiseux-first-coefficients}.  For positive real
$z$ sufficiently small, the plus branch has imaginary part strictly
less than $\theta$ and real part inside $(0,M)$; it lies in $D_M$.
The minus branch lies above the upper horizontal side.  Since
$|\tau(1-z)|<R$, uniqueness inside $D_M$ identifies the plus branch
with $U_A$.  The coefficients $b,c_3$ follow by substituting the
Puiseux series into \eqref{eq:local-implicit-equation}.  Complex
conjugation settles the lower endpoint.  The nonzero square-root
coefficient makes both endpoints genuine singularities, and the
radius is exactly $R$.

For completeness, the local continuations give the full slit domain
in \eqref{eq:slit-domain}.  Take disjoint small disks around the two
critical endpoints, with the selected Puiseux branches cut along the
outward radial rays.  At each regular point on $|t|=R$, choose an
implicit-function disk centered there and small enough to avoid both
critical endpoints.  Such a disk avoids both radial cuts: the point
of an outward ray closest to its center is the corresponding endpoint.
Two overlapping disks whose centers lie on $|t|=R$ have an overlap
meeting $|t|<R$, as seen by taking a suitable point of the chord between
their centers.  Regular overlaps, including regular/critical ones,
are connected and contain no cut, while the critical disks are disjoint.
Thus all analytic branches agree on overlaps by the identity theorem.
A finite cover of the circle by the underlying uncut disks has a
positive collar width.  Subtracting the radial rays from this collar
proves \eqref{eq:slit-domain} after decreasing $\epsilon$.
\end{proof}

The proof selects a bounded inverse sheet directly.  In particular, it
does not minimize the moduli of all formal critical values
$-A/W_k(-A/e)$, which accumulate at zero on other sheets.  That
minimization would be invalid for the germ under consideration.  The
radius agrees with \cite[Eq.~(2.9)]{KaluginJeffrey2012} under
\eqref{eq:classical-variable-map}; the barrier supplies the additional
continuation and branch information used below.

\section{Unconditional oscillatory asymptotics to arbitrary order}
\label{sec:velocity-asymptotics}

\begin{theorem}[Two terms and all orders]
\label{thm:velocity-asymptotics}
Fix $A>1$ and let the parameters be as in
\eqref{eq:critical-parameters}--\eqref{eq:puiseux-first-coefficients}.
Define
\begin{equation}\label{eq:chi-correction}
 \chi_1=\chi\left(-\frac38+\frac{u_*}{12}+\frac3{8u_*}\right).
\end{equation}
Then
\begin{equation}\label{eq:velocity-two-term}
 v_n(A)=\frac{R^{-n}}{\sqrt\pi\,n^{3/2}}
    \Re\left[e^{-in\theta}
                      \left(\chi+\frac{\chi_1}{n}\right)\right]
           +O_A(R^{-n}n^{-7/2}).
\end{equation}
In particular, with $\phi=-\arg\chi\in(\pi/4,\pi/2)$,
\begin{equation}\label{eq:velocity-leading-cosine}
 v_n(A)=\frac{|\chi|}{\sqrt\pi\,R^n n^{3/2}}
               \cos(n\theta+\phi)+O_A(R^{-n}n^{-5/2}).
\end{equation}
These are additive estimates, including at indices where the displayed
cosine is small.  An expansion to any prescribed number of inverse
powers of $n$ follows from explicitly computable local coefficients.
\end{theorem}

\begin{proof}
At $\tau$, the nonintegral terms of $V_A=L-U_A$ are
$-\chi(1-t/\tau)^{1/2}-c_3(1-t/\tau)^{3/2}$, with conjugate terms
at $\overline\tau$.  The slit domain from
Theorem~\ref{thm:selected-sheet} permits the standard finite-singularity
transfer theorem \cite[Chapter~VI]{FlajoletSedgewick2009}.  For a fixed
nonintegral exponent $\alpha$,
\[
 [t^n](1-t/\tau)^\alpha
 =\frac{\tau^{-n}n^{-\alpha-1}}{\Gamma(-\alpha)}
     \left(1+\frac{\alpha(\alpha+1)}{2n}+O(n^{-2})\right).
\]
Using $\Gamma(-1/2)=-2\sqrt\pi$ and
$\Gamma(-3/2)=4\sqrt\pi/3$, the upper singularity contributes
\[
 \frac{\tau^{-n}}{2\sqrt\pi\,n^{3/2}}
       \left(\chi+\frac{3\chi/8-3c_3/2}{n}\right)
       +O_A(R^{-n}n^{-7/2}).
\]
Adding its conjugate and substituting for $c_3$ gives
\eqref{eq:velocity-two-term}.

For an explicit all-orders construction, put
\[
 K(w)=\frac{2u_*(e^w-1-w)}{w^2(u_*+w)},\qquad K(0)=1.
\]
Equation \eqref{eq:local-implicit-equation} becomes
$w=\chi z^{1/2}K(w)^{-1/2}$.  Lagrange inversion gives the convergent
Puiseux coefficients
\begin{equation}\label{eq:all-puiseux-coefficients}
 U_A(t)=u_*+\sum_{m\geq1}c_mz^{m/2},\qquad
 c_m=\frac{\chi^m}{m}[w^{m-1}]K(w)^{-m/2}.
\end{equation}
For any integer $J\geq1$, transfer of the odd coefficients yields the
particularly convenient form
\begin{equation}\label{eq:binomial-all-order}
 v_n(A)=-2\Re\left[\tau^{-n}
       \sum_{j=0}^{J-1}c_{2j+1}(-1)^n\binom{j+\tfrac12}{n}\right]
       +O_{A,J}(R^{-n}n^{-J-3/2}).
\end{equation}
The next omitted nonintegral exponent is $J+1/2$; the integer powers
are locally analytic and do not contribute algebraic transfer terms.
Expanding the exact binomial coefficients by gamma-ratio asymptotics
gives the requested expansion in inverse powers of $n$.
\end{proof}

For $A=2$, the exact angle definition gives numerically
\[
 u_*=0.469353635658274\ldots+1.132672497604881\ldots i,
 \qquad R=1.598960348180173\ldots,
\]
and
\[
 \chi=0.869892622200723\ldots-1.302083117729355\ldots i.
\]
The two-term formula in \cite{HarmonicBaseline} is therefore now
unconditional.  These decimals merely illustrate the exact parameters;
they have no role in the continuation proof.

\section{Bounded sign gaps and a complete density alternative}
\label{sec:velocity-signs}

The earlier infinite-sign theorem leaves open the lengths of possible
runs in one direction.  The singularity analysis gives a uniform bound
on such runs for each fixed lattice spacing.

\begin{corollary}[Every sufficiently late short block has both signs]
\label{cor:bounded-sign-gaps}
For $A>1$ put
\[
 K_A=\left\lfloor\frac\pi\theta\right\rfloor+2.
\]
Every sufficiently late block of $K_A$ consecutive coefficients $v_n(A)$
contains a strictly positive and a strictly negative coefficient.  In
particular, every sufficiently late block of four coefficients $v_n(2)$
has both signs.
\end{corollary}

\begin{proof}
Put $m=K_A-1$.  Then $\pi<m\theta\leq\pi+\theta<2\pi$.
The points $0,\theta,\ldots,m\theta$ on the circle have largest
circular gap
\[
 g=\max\{\theta,2\pi-m\theta\}<\pi.
\]
Choose $\eta>0$ with $\pi-2\eta>g$.  Every rotation of these points
meets both an arc where $\cos x\geq\sin\eta$ and an arc where
$\cos x\leq-\sin\eta$, since each arc has length $\pi-2\eta$.
After division by the positive prefactor in
\eqref{eq:velocity-leading-cosine}, the error is $O_A(1/n)$ and is
eventually smaller than $\sin\eta/2$.  Thus both strict signs occur.

For the final assertion, $\pi/3<\theta(2)<\pi/2$ can be proved without
decimal computation.  At $\theta=\pi/3$ put
$x=\pi/(3\sqrt3)\in(0,1)$.  The corresponding value of $A$ is
$2xe^{1-x}<2$, by $\log x<x-1$.  At $\theta=\pi/2$ it is
$\pi e/2>2$.  Monotonicity in \eqref{eq:theta-definition} gives the
claimed bracket and hence $K_2=4$.
\end{proof}

\begin{theorem}[Sign densities for every lattice parameter]
\label{thm:sign-density}
If $\theta/(2\pi)$ is irrational, the positive and negative velocities
each have natural density $1/2$, and the zero velocities have density
zero.  If $\theta/(2\pi)=p/q$ is rational in lowest terms, the velocity
signs are eventually periodic of period dividing $q$, none is eventually
zero, and their densities are
\begin{equation}\label{eq:rational-sign-density}
 d_+(A)=\frac1q\#\{0\leq j<q:
                     \cos(2\pi j/q+\phi)>0\},\qquad
 d_-(A)=1-d_+(A).
\end{equation}
For even $q$ both densities are $1/2$; for odd $q$ the positive count is
$(q-1)/2$ or $(q+1)/2$.
\end{theorem}

\begin{proof}
In the irrational case the sequence $n\theta+\phi$ is equidistributed
modulo $2\pi$ by Weyl's theorem \cite{Weyl1916}.  For any $\delta>0$, the additive estimate
\eqref{eq:velocity-leading-cosine} implies that any disagreement with
the cosine sign, including a zero coefficient, eventually occurs only
where $|\cos(n\theta+\phi)|<\delta$.  The density of this set tends to
zero with $\delta$.  The asserted densities follow.

In the rational case it is essential to rule out a zero among the finite
leading cosine values.  Such a zero would make $\phi/\pi$ rational.
Since $\arg u_*=\pi-2\phi$, the positive number
$b_0=\theta/a=\tan(\arg u_*)$ would then be algebraic.  Also
$\cot\theta$ is algebraic when $\theta/\pi$ is rational.  But
$a=1-\theta\cot\theta$ implies
\[
 \theta(1+b_0\cot\theta)=b_0.
\]
The coefficient of $\theta$ cannot vanish, because $b_0>0$.  This would
make $\theta$ algebraic, contradicting the transcendence of $\pi$ \cite{Lindemann1882}.
Thus every leading cosine has a nonzero margin from zero, and the
$O_A(1/n)$ error gives eventual periodicity.  Multiplication by $p$
permutes the $q$ phase points, proving \eqref{eq:rational-sign-density}.
Antipodal pairing gives the even-$q$ statement.  For odd $q$, an open
semicircle avoiding the sample points contains $(q-1)/2$ or $(q+1)/2$
of the equally spaced points.
\end{proof}

The rational-angle parameters form the explicit countable set
\[
 A=\frac\theta{\sin\theta}\exp(1-\theta\cot\theta),
 \qquad \theta=2\pi p/q\in(0,\pi).
\]
No irrationality claim is made for the particular angle associated with
$A=2$.  The four-term sign-block result is unconditional and does not
require deciding that arithmetic question.

\section{Convergent identities at regular and critical parameters}
\label{sec:convergent-velocity-identities}

The preceding estimates turn formal substitutions in the inverse
equation into justified convergent identities.  They also describe the
first differentiated series at the critical endpoints.

\begin{proposition}[Interior sums and critical boundary values]
\label{prop:velocity-summations}
For every $A>1$, both of the following series converge absolutely:
\begin{align}
 \sum_{n\geq1}v_n(A)
    &=\log A+A+W_{-1}(-e^{-A}),
                       \label{eq:sum-plus-one}\\
 \sum_{n\geq1}(-1)^n v_n(A)
    &=\log A-A+W_0(e^A).
                       \label{eq:sum-minus-one}
\end{align}
On the critical circle the series still converges absolutely, and
\begin{align}
 \sum_{n\geq1}v_n(A)R^n\cos(n\theta)&=\log A-a,
                           \label{eq:critical-cosine-sum}\\
 \sum_{n\geq1}v_n(A)R^n\sin(n\theta)&=-\theta.
                           \label{eq:critical-sine-sum}
\end{align}
The differentiated critical series has the precise finite-part limit
\begin{equation}\label{eq:critical-finite-part}
 \sum_{n=1}^N n v_n(A)\tau^n
   =\frac{\chi}{\sqrt\pi}\sqrt N+\frac{u_*}{3}-1
                  +O_A(N^{-1/2}).
\end{equation}
\end{proposition}

\begin{proof}
Since $a>0$, we have $R>1$, so substitution of $t=\pm1$ lies strictly
inside the disk of convergence.  On the positive real $U$ axis the
solutions at those parameters are
\[
 U_A(1)=-A-W_{-1}(-e^{-A}),\qquad
 U_A(-1)=A-W_0(e^A).
\]
They are positive and satisfy the implicit equation.  They are the
selected roots by uniqueness in the rectangle, or by continuation along
the real interval, where $U_A>0$.  Substitution into $V_A=\log A-U_A$
proves \eqref{eq:sum-plus-one}--\eqref{eq:sum-minus-one}.

Equation \eqref{eq:velocity-leading-cosine} gives
$\sum_{n\geq1}|v_n(A)|R^n<\infty$.  Abel continuity at $\tau$ and
Theorem~\ref{thm:selected-sheet} therefore yield
$\sum_{n\geq1}v_n(A)\tau^n=\log A-u_*$.  Real and imaginary parts
give \eqref{eq:critical-cosine-sum}--\eqref{eq:critical-sine-sum}.

Finally write $r=t/\tau$ and
$B(r)=r\,dV_A(\tau r)/dr=\sum_{n\geq1}nv_n(A)\tau^n r^n$.
The expansion at $r=1$ is
\[
 B(r)=\frac\chi2(1-r)^{-1/2}+b+O((1-r)^{1/2}).
\]
The partial sum in \eqref{eq:critical-finite-part} is
$[r^N]B(r)/(1-r)$.  Its singular terms at $r=1$ give
\[
 \frac\chi2\,[r^N](1-r)^{-3/2}+b
     =\frac\chi{\sqrt\pi}\sqrt N+b+O_A(N^{-1/2}).
\]
The other singularity is at $r=\overline\tau/\tau=e^{-2i\theta}\neq1$.
There the factor $1/(1-r)$ is regular and $B$ has exponent $-1/2$,
so its coefficient contribution is $O_A(N^{-1/2})$.  Transfer of the
remaining terms has the same or smaller order.  This proves the
finite-part formula.
\end{proof}

All sums in this proposition become explicit Stirling identities on
using \eqref{eq:velocity-stirling}.  Their derivation is a consequence
of the proved continuation and classical inverse representation;
no claim of independent priority for each specialized identity is made.
